\bookchapter{פרק 39: הנדסת הקשר — לגרום ל-AI להבין את כוונתכם}

\chapterintro{האם קרה לכם: פתחתם חלון שיחה חדש, ופתאום ה-AI נראה כמי שסבל מאמנזיה ושכח לגמרי את כל מה שקשור לפרויקט שלכם? זה לא בגלל שה-AI מנסה להרגיז אתכם — זו פשוט “שכחה” מולדת שלו. בפרק זה תלמדו להתקין ל-AI “מערכת זיכרון ארוכת טווח”, ובדרך תכירו אחת מפרדיגמות הפיתוח החמות של 2026 — פיתוח מונחה מפרט (Spec-Driven Development).}

\booksection{א. קובץ זיכרון הפרויקט: “הזיכרון ארוך הטווח” של ה-AI}

למודלי AI יש מגבלה יסודית: הזיכרון שלהם “קצר”. בכל תחילת שיחה הידע של ה-AI על הפרויקט שלכם שווה אפס. עליכם להסביר לו שוב ושוב בתוך השיחה “אנחנו עובדים עם React”, “הבק-אנד הוא Node.js”, “מסד הנתונים הוא PostgreSQL”, “סגנון הקוד — מירכאות כפולות”. זה לא רק בזבזני, אלא גם מזמן השמטות — כמו לערוך מדי יום קבלת עובדים מחדש לעמית חדש.

הפתרון הוא “קובץ זיכרון הפרויקט” — קובץ טקסט פשוט בתיקיית השורש של הפרויקט, שכלי ה-AI קוראים אוטומטית בכל תחילת שיחה. לכל כלי יש מוסכמת שמות משלו:

\begin{center}
\begin{tabularx}{\linewidth}{|X|X|}
\hline
\textbf{כלי} & \textbf{שם הקובץ} \tabularnewline
\hline
Cursor & \texttt{.\allowbreak cursorrules} \tabularnewline
Claude Code & \texttt{CLAUDE.\allowbreak md} \tabularnewline
GitHub Copilot & \texttt{copilot-\allowbreak instructions.\allowbreak md} \tabularnewline
פתרון כללי & \texttt{AGENTS.\allowbreak md} \tabularnewline
\hline
\end{tabularx}
\end{center}

\begin{center}\small\color{gray}מוסכמות השמות של קובצי זיכרון הפרויקט בכלי השונים\end{center}

אמנם השמות שונים, אבל העיקרון אחד: את המידע שאתם רוצים שה-AI “יזכור תמיד” כותבים בקובץ הזה. קובץ זיכרון פרויקט טוב אמור לכלול: סקירת הפרויקט (מהו הפרויקט), מחסנית טכנולוגיות (במה משתמשים), כללי קוד (שמות, הזחה, סגנון הערות), מבנה תיקיות (איך הקוד מאורגן) ומוסכמות פיתוח (אסטרטגיית ענפים, כללי קומיט, דרישות בדיקות). ה-CLAUDE.md של Karpathy מכיל 70 שורות בלבד, ומכסה את כל המידע הקריטי — תמציתי אך ממצה.

\begin{lstlisting}
# פרויקט: אפליקציה לניהול הוצאות אישיות
## סקירה
אפליקציית full-stack לניהול הוצאות אישיות שעוזרת למשתמש לתעד הכנסות והוצאות יומיות.
## סטאק טכנולוגי
- Frontend: React + TypeScript + Tailwind CSS
- Backend: Node.js + Express + TypeScript
- מסד נתונים: SQLite (פיתוח) / PostgreSQL (ייצור)
- בדיקות: Vitest
## כללי קוד
- להשתמש במרכאות כפולות
- הזחה של 2 רווחים
- שמות פונקציות ב-camelCase
- הגדרות טיפוסים ב-PascalCase
- לכל פונקציה חייב להיות תיעוד JSDoc
## מבנה תיקיות
/src
  /frontend  - קוד הפרונט
  /backend   - קוד הבק
  /shared    - טיפוסים משותפים לפרונט ולבק
  /tests     - קבצי בדיקה
## הסכמי פיתוח
- הודעות commit בעברית
- כל יכולת חדשה מתחילה בבדיקה
- תשובות ה-API תמיד בפורמט { success, data, error }
\end{lstlisting}

\begin{tcbwarn}{שימו לב: אל תהפכו אותו ל"מסמך דרישות שני"}
שימו לב לאורך הקובץ — אל תכתבו מסמך דרישות בן עשרות עמודים. מטרת קובץ זיכרון הפרויקט היא “לאפשר ל-AI לבנות במהירות תמונת רקע”, לא “לתאר כל פיצ'ר לעומק”. שמרו על תמציתיות, על מבנה ברור, וכתבו רק את “מה שה-AI צריך לדעת”.

\end{tcbwarn}

\booksection{ב. פיתוח מונחה מפרט (SDD)}

Spec-Driven Development (SDD, פיתוח מונחה מפרט) היא פרדיגמת פיתוח שהבשילה בהדרגה ב-2026, ורעיון הליבה שלה במשפט אחד: קודם כותבים מפרט, ואז כותבים קוד. בשונה מ-TDD המסורתי (פיתוח מונחה בדיקות), ה“מפרט” ב-SDD אינו רק מקרי בדיקה, אלא כולל גם תיאור דרישות, הגדרות ממשק, מודלי נתונים ועוד. ה-AI מחולל קוד על סמך המפרטים האלה, ותפקיד בני האדם הוא לכתוב ולתחזק את המפרט.

תהליך העבודה של SDD בחמישה צעדים:

\begin{enumerate}
  \item מתארים את מפרט הדרישות בשפה טבעית (בתוך קובץ .md);
  \item מגדירים ממשקים ומבני נתונים (בתוך קובץ המפרט);
  \item מבקשים מה-AI לחולל את מימוש הקוד על סמך המפרט;
  \item מריצים בדיקות לאמת שהתוצר תואם למפרט;
  \item אם אין התאמה — מתקנים את המפרט או מבקשים מה-AI לתקן את הקוד.
\end{enumerate}

בתהליך הזה קובץ המפרט הוא “מקור האמת היחיד” — את הקוד אפשר לחולל מחדש, אבל המפרט חייב להיות נכון תחילה.

כתיבת מסמך המפרט

האם המפרט מלא?

מלא

ה-AI מייצר קוד על פי המפרט

להריץ בדיקות כדי לוודא את התוצר

לא מלא / הבדיקות נכשלות

בעיה בקוד תקן את הקוד; בעיה במפרט תקן את המפרט

חזרה לבדיקה

הכול עבר מסירה

\begin{center}\small\color{gray}איור 3: תהליך העבודה של פיתוח מונחה מפרט (SDD)\end{center}

יתרונו של SDD: קובץ המפרט יציב יותר מהקוד. כשהמודל או הכלי משתנים, קובץ המפרט שלכם נותר בתוקף, וכל שנדרש הוא לבקש מכלי AI חדש לחולל את הקוד מחדש. וזו בדיוק הסיבה שהספר הזה מדגיש “המתודולוגיה חשובה יותר מהכלים” — קובץ המפרט הוא הנשא המוחשי של המתודולוגיה שלכם.

\booksection{ג. ניהול חלון ההקשר}

גם עם קובץ זיכרון פרויקט, עדיין תיתקלו בשיחות בבעיית “אין מספיק הקשר”. חלון ההקשר של מודל AI (כמות המידע שהוא מסוגל לעבד בו-זמנית) הוא מוגבל. כשהשיחה ארוכה מדי ומוזכרים בה יותר מדי קבצים, ה-AI “שוכח” מידע מוקדם ואיכות התוצר יורדת — כמו חבר שאיתו שוחחתם שעות ושכבר שכח את הדבר הראשון שסיפרתם לו בתחילת היום.

האסטרטגיה המרכזית לניהול חלון ההקשר היא “התייחסות בשכבות”:

\begin{center}
\begin{tabularx}{\linewidth}{|X|X|X|}
\hline
\textbf{שכבת ההקשר} & \textbf{תוכן} & \textbf{אופן הניהול} \tabularnewline
\hline
שכבה ראשונה: זיכרון הפרויקט & מחסנית טכנולוגיות, מפרטים, מבנה תיקיות & נכתב בקבצים כמו CLAUDE.md, נטען אוטומטית ואין צורך לחזור עליו \tabularnewline
שכבה שנייה: התייחסות לקבצים & קובצי הקוד הנוגעים למשימה הנוכחית & מפנים במדויק באמצעות סימן @, בלי להדביק את הקובץ המלא לשיחה \tabularnewline
שכבה שלישית: הקשר השיחה & היסטוריית השיחה של המשימה הנוכחית & אחרי השלמת משימה פותחים שיחה חדשה, כדי למנוע נפיחות הקשר \tabularnewline
שכבה רביעית: הוראות מיידיות & הדרישות הספציפיות של הצעד הנוכחי & בהירות, תמציתיות ומיקוד, בלי תיאורים ארוכים \tabularnewline
\hline
\end{tabularx}
\end{center}

\begin{center}\small\color{gray}אסטרטגיית הניהול השכבתי של חלון ההקשר\end{center}

\booksection{ד. נקודות מפתח בתרגול של הנדסת הקשר}

הנדסת הקשר אינה הגדרה חד-פעמית אלא תהליך של שיפור מתמיד. הנה כמה נקודות שאומתו שוב ושוב בשטח:

\begin{enumerate}
  \item מעדכנים את קובץ הזיכרון בהתאמה עם התפתחות הפרויקט. כשמאמצים טכנולוגיה חדשה או משנים ארכיטקטורה — מעדכנים מיד את קובץ הזיכרון, במקום להסביר זאת באופן חד-פעמי בתוך השיחה.
  \item מנצלים התייחסות ב-@ במקום הדבקת טקסט מלא. רוב כלי התכנות בעזרת AI תומכים בתחביר @file או @folder, שמזריקים את תוכן הקובץ להקשר של ה-AI ברקע — יעיל יותר מהדבקה ידנית ולא מזהם את השיחה.
  \item שומרים על שיחות "קצרות וממוקדות". שיחה אחת מטפלת במשימה אחת או בנקודת פונקציונליות אחת. בסיומה פותחים שיחה חדשה למשימה הבאה. שיחה ארוכה לא רק שוחקת את חלון ההקשר, אלא גם מסיחה את דעת ה-AI.
  \item נותנים ל-AI "לראות" את התוצאה ולא "לשמוע תיאור" שלה. כשנתקלים בבאג, מפנים את יומן השגיאות ישירות ב-@ ל-AI, במקום לתאר במילים שלכם "נראה לי שיש איזו שגיאה". ל-AI קל יותר לעבד את המידע הגולמי מאשר את תמציתכם.
  \item מנקים ומארכיבים באופן קבוע. שיחות שהסתיימו וקבצים זמניים תופסים מקום עבודה. ניקוי קבוע של תוכן לא רלוונטי ושמירה על תיקיית פרויקט מסודרת עוזרים גם ל-AI להבין את מבנה הפרויקט ביעילות רבה יותר.
\end{enumerate}

\begin{tcbpractice}{תרגול מעשי: כתבו לפרויקט שלכם קובץ זיכרון בן 70 שורות}
בחרו פרויקט שאתם עובדים עליו (או מתכננים להתחיל), וכתבו קובץ CLAUDE.md: סקירת הפרויקט (2-3 משפטים), מחסנית הטכנולוגיות (רשימה), כללי קוד (3-5 סעיפים), מבנה תיקיות (אילן קטן) ומוסכמות פיתוח (2-3 סעיפים). הקפידו לא לחרוג מ-70 שורות — בסיום תגלו שגם ההבנה שלכם את הפרויקט התבהרה משמעותית.

\end{tcbpractice}

\begin{tcbsummary}{סיכום הפרק}
\begin{itemize}
  \item הזיכרון של ה-AI הוא "קצר": בכל שיחה חדשה הידע שלו על הפרויקט שווה אפס
  \item קובץ זיכרון הפרויקט (CLAUDE.md ודומיו) הוא "הזיכרון ארוך הטווח" של ה-AI, ונטען אוטומטית בכל שיחה
  \item קובץ זיכרון טוב = סקירה + מחסנית טכנולוגיות + מפרטים + תיקיות + מוסכמות — 70 שורות מספיקות
  \item SDD — פיתוח מונחה מפרט: קודם מפרט ואז קוד; המפרט הוא "מקור האמת היחיד" ויציב יותר מהקוד
  \item ההקשר מנוהל בארבע שכבות: זיכרון הפרויקט התייחסות לקבצים ב-@ היסטוריית שיחה הוראות מיידיות
  \item חמש נקודות תרגול: עדכון מסונכרן, התייחסות ב-@, שיחות קצרות, מתן מידע גולמי, ניקוי תקופתי
\end{itemize}

\end{tcbsummary}
